\section{Temporal Circuit Hierarchies and the Distinction Between Deliberation and Execution}
\label{sec:temporal-hierarchies}

% ═══════════════════════════════════════════════════════════════════
\subsection{Multiple Nested Circuit Closures}
% ═══════════════════════════════════════════════════════════════════

A closed microfluidic network does not close at a single timescale. Instead, it admits multiple nested circuit closures operating at distinct temporal frequencies, each characterized by different latency requirements and topological constraints.

\begin{definition}[Temporal closure hierarchy]
\label{def:temporal-hierarchy}
A closed microfluidic network exhibits a temporal hierarchy of closures:

\begin{enumerate}[label=(\roman*),leftmargin=*,itemsep=2pt]

\item \textbf{Fast reflex closure} ($\tau_{\text{fast}} \sim 30\text{--}100\,\text{ms}$):
Local sensorimotor loops operating at the spinal and brainstem level. Motor command → muscle → proprioceptor → sensory afferent → motor nucleus → corrective command. These loops close without requiring input from higher cognitive centers. Example: stretch reflex arc.

\item \textbf{Medium coordination closure} ($\tau_{\text{med}} \sim 0.2\text{--}1\,\text{s}$):
Spinal interneuron networks and cerebellar loops that synchronize activity across multiple reflex arcs. These operate at intermediate timescale, coordinating the outputs of many fast loops through shared descending and ascending pathways.

\item \textbf{Slow deliberative closure} ($\tau_{\text{slow}} \sim 1\text{--}5\,\text{s}$):
Supraspinal circuits (motor cortex, prefrontal cortex, cerebellum, vestibular nuclei, brainstem) integrating multimodal sensory input, internal state, and goal representations. This closure includes working memory, planning, prediction, and the issuance of descending motor commands that bias the operating point of the faster closures.

\end{enumerate}

Each closure is topologically a circuit: charge redistribution at that level must return to its origin. Each operates at its characteristic timescale and is driven by metabolic energy input.
\end{definition}

\begin{theorem}[Timescale separation prevents shared execution]
\label{thm:timescale-separation}
Let $S_{\text{fast}}$ be a computational substrate executing fast reflex closure ($\tau_{\text{fast}} \le 100\,\text{ms}$) and let $S_{\text{slow}}$ be a substrate executing slow deliberative closure ($\tau_{\text{slow}} \ge 1\,\text{s}$). If both closures operate on the same critical path (same sensorimotor output), then no single substrate can satisfy the latency requirements of both simultaneously. Therefore, the two closures must operate on distinct but coupled pathways.
\end{theorem}

\begin{proof}
Suppose both closures shared the same effector pathway. Then the system's output latency would be dominated by whichever closure is slower. If the slow closure is on the critical path, output latency $\tau_{\text{out}} \ge \tau_{\text{slow}} \ge 1\,\text{s}$. But the fast closure requires $\tau_{\text{out}} \le 100\,\text{ms}$ to stabilize high-bandwidth disturbances. These constraints are incompatible; hence the slow closure cannot be on the fast closure's critical path.

The only resolution is for the slow closure to operate off-critical-path, modulating parameters (e.g., bias, gain, set-point) that the fast closure consumes asynchronously. The fast closure closes on its own timescale; the slow closure provides time-varying boundary conditions.
\end{proof}

\begin{corollary}[Hierarchical coupling]
The fast closure must be capable of closing without real-time input from the slow closure. The slow closure modifies the steady-state operating point of the fast closure, but does not gate its execution.
\end{corollary}

\subsection{Charge Redistribution at Different Timescales}

Each closure involves charge redistribution at its characteristic frequency.

\begin{definition}[Fast-closure charge rate]
Charge redistribution in fast reflex loops (stretch reflex, cutaneous reflex, Golgi tendon reflex) occurs at frequencies $f_{\text{fast}} \sim 10\text{--}30\,\text{Hz}$ with each loop completing in 30--100 ms. The charge per reflex cycle is modest (single motor units, local membrane depolarization):
\begin{equation}
Q_{\text{fast,cycle}} \sim 1\text{--}10\,\text{nC per reflex arc}
\end{equation}
\end{definition}

\begin{definition}[Medium-closure charge rate]
Charge redistribution in spinal coordination networks operates at $f_{\text{med}} \sim 1\text{--}5\,\text{Hz}$ (multiple reflex arcs synchronized). The charge per cycle aggregates across many spinal segments:
\begin{equation}
Q_{\text{med,cycle}} \sim 100\,\text{nC per coordination cycle}
\end{equation}
\end{definition}

\begin{definition}[Slow-closure charge rate]
Charge redistribution in supraspinal deliberative circuits (motor cortex, prefrontal, cerebellum, vestibular nuclei) operates at $f_{\text{slow}} \sim 0.2\text{--}2\,\text{Hz}$. The charge per cycle is large, involving millions of cortical neurons:
\begin{equation}
Q_{\text{slow,cycle}} \sim 100\text{--}300\,\text{mC per deliberative cycle}
\end{equation}

This is the charge rate quantified in Section 6: $Q_{\text{thought}} = 133.5\,\text{mC/s}$ corresponds to approximately one deliberative cycle per second.
\end{definition}

% ═══════════════════════════════════════════════════════════════════
\subsection{Hierarchical Topological Coupling}
% ═══════════════════════════════════════════════════════════════════

\begin{theorem}[Nested closure with parametric modulation]
\label{thm:nested-closure}
Let $\mathcal{C}_{\text{fast}}(t)$ denote the fast closure trajectory, parametrized by a bias signal $b(t)$ from the slow closure: $\dot{\mathcal{C}}_{\text{fast}} = F_{\text{fast}}(\mathcal{C}_{\text{fast}}, b(t))$. Let $\mathcal{C}_{\text{slow}}(t)$ denote the slow closure trajectory. Then:

\begin{enumerate}[label=(\alph*),leftmargin=*,itemsep=1pt]
\item The fast closure admits a family of steady-state limit cycles parametrized by $b$: $\{\mathcal{C}_{\text{fast}}^{\star}(b) : b \in B\}$.
\item The slow closure generates $b(t)$ via integration of supraspinal signals, sampling the attractor basin of steady-state limit cycles.
\item The fast closure's execution is independent of $b(t)$'s rate of change (provided $b$ changes slowly relative to $\tau_{\text{fast}}$).
\item Feedback from the fast closure (proprioceptive, vestibular, visual) enters the slow closure as a time-averaged input, not cycle-by-cycle.
\end{enumerate}
\end{theorem}

\begin{proof}
(a)--(b): By the timescale-separation theorem, if $\tau_{\text{slow}} \gg \tau_{\text{fast}}$, the slow dynamics act as a quasi-static driver of parameters in the fast system. For each fixed $b$, the fast closure converges to a limit cycle $\mathcal{C}_{\text{fast}}^{\star}(b)$ within time $\sim \tau_{\text{fast}}$. As $b$ varies on the slow timescale, the system tracks a family of limit cycles.

(c): The fast closure is driven by $F_{\text{fast}}(\mathcal{C}_{\text{fast}}, b)$. Within one fast cycle (duration $\tau_{\text{fast}} \sim 0.1\,\text{s}$), $b(t)$ changes by $\Delta b \sim \dot{b} \cdot \tau_{\text{fast}} \ll b$ (since $\dot{b}$ is set by slow dynamics). The fast cycle executes with $b$ approximately constant; the limit cycle is insensitive to the slow change.

(d): Proprioceptive and vestibular feedback provide dense sampling of the body state, but the slow closure integrates this feedback on the slow timescale, not instantaneously. The return signal to the slow loop is a smoothed estimate $\bar{s}(t) = \int_t^{t+\tau_{\text{slow}}} s(\tau)\,d\tau / \tau_{\text{slow}}$, not the raw fast-loop state.
\end{proof}

\noindent\textbf{Interpretation:} The slow closure does not directly control the fast closure. Instead, it sets the operating point (bias, reference configuration, motor set). The fast closure then executes that point robustly, adjusting moment-to-moment via proprioceptive feedback. Once the slow closure has issued a bias signal, the fast loops run to completion without further revision.

% ═══════════════════════════════════════════════════════════════════
\subsection{Deliberation Versus Execution}
% ═══════════════════════════════════════════════════════════════════

The distinction between what is commonly called ``thought'' and ``action'' maps precisely onto the distinction between slow and fast closures.

\begin{definition}[Deliberative charge redistribution]
Deliberative charge redistribution is the slow-closure activity ($\tau_{\text{slow}} \sim 1\text{--}5\,\text{s}$) in which the supraspinal network:
\begin{enumerate}[label=(\roman*),leftmargin=*,itemsep=1pt]
\item Integrates multimodal sensory information (visual, vestibular, proprioceptive, interoceptive, exteroceptive).
\item Maintains internal state representation (goals, world model, body model, history).
\item Explores multiple possible motor commands and predicts their outcomes.
\item Selects a bias signal $b^{\star}$ that the system commits to execute.
\end{enumerate}
This is the process popularly termed ``thinking'' or ``decision-making.''
\end{definition}

\begin{definition}[Executive charge redistribution]
Executive charge redistribution is the fast-closure activity ($\tau_{\text{fast}} \sim 30\text{--}100\,\text{ms}$) in which the spinal and brainstem networks:
\begin{enumerate}[label=(\roman*),leftmargin=*,itemsep=1pt]
\item Receive the bias signal $b(t)$ from the slow closure.
\item Close sensorimotor loops against that bias: generate muscle commands, receive proprioceptive feedback, correct deviations.
\item Execute the commanded action with sub-second latency.
\end{enumerate}
This is the process popularly termed ``action'' or ``motor execution.''
\end{definition}

\begin{remark}[The F1 driver paradigm]
A Formula~1 driver at the braking point does not decide ``turn the steering wheel 0.5 degrees now'' every 50 ms. Instead:
\begin{itemize}
\item The slow closure (1--5 s) decides the overall strategy: ``approach this corner at speed $v$, apex at location $(x,y)$, exit at velocity $v'$.''
\item This decision is transmitted as bias signals to the motor system.
\item The fast closure (30--100 ms) executes: steering wheel angle adjusts in real time via proprioceptive feedback, balancing the car against the tire forces, correcting for wind gusts and track variations.
\item The driver's conscious experience is the slow-closure deliberation. The steering adjustments happen below consciousness.
\end{itemize}
\end{remark}

% ═══════════════════════════════════════════════════════════════════
\subsection{Production-Completion Incompatibility at Nested Timescales}
% ═══════════════════════════════════════════════════════════════════

From Section 5, we know that within a single closure, knowledge production ($\sigma > 0$: exploring possibilities) and action commitment ($\sigma = 0$: executing a selected action) are mutually exclusive at each iteration. This principle extends across the hierarchy.

\begin{theorem}[Hierarchical production-completion]
\label{thm:hierarchical-incompatibility}
At the slow-closure timescale, the supraspinal network cannot simultaneously deliberate (explore multiple possible bias signals $b_1, b_2, \ldots$) and issue a committed bias signal. Once the slow closure commits to $b^{\star}$, the deliberative state becomes fixed at that value for the duration of fast-loop execution. Revision requires suspending fast execution and re-entering deliberation.
\end{theorem}

\begin{proof}
Within the slow-closure timescale $\tau_{\text{slow}} \sim 1\text{--}5\,\text{s}$, deliberation involves activation of multiple competing motor representations in prefrontal cortex, premotor areas, and parietal cortex. These represent the multiple possible bias signals. Issuing a committed bias signal requires selection of one representation and silencing of the others, transmitted via motor cortex and brainstem descending pathways.

The commitment is transmitted to the fast closure as $b^{\star}(t)$, which the fast loops now track. If the slow closure were to simultaneously issue competing bias signals (high deliberative variance), the fast loops would receive conflicting inputs and could not close stably. Therefore, at the moment of commitment, deliberative variance must drop to zero, and $b^{\star}$ must be held constant for the duration of the action.

Revising the commitment (changing $b^{\star}$ before the fast action completes) requires:
1. Suppressing the ongoing fast-loop execution (motor command abort).
2. Re-entering the deliberative state at the slow closure (variance rises again).
3. Exploring new bias candidates.
4. Recommitting to a new $b^{\star\star}$.

These steps cannot occur simultaneously.
\end{proof}

\subsection{The Intentionality Principle: Harmful Outcomes Reveal Circuit Failure}

\begin{principle}[Harmful outcomes reveal circuit failure, not intention]
\label{prin:intentionality}
The supraspinal slow closure only commits to a bias signal $b^{\star}$ if the associated prediction is:
\begin{enumerate}[label=(\roman*),leftmargin=*,itemsep=1pt]
\item Consistent with stated goals (reaching a target, maintaining balance, avoiding obstacles).
\item Predicted to execute without unintended harm.
\end{enumerate}

If an action results in unintended harm (falling, crashing, injury), it reveals one of three possibilities:
\begin{enumerate}[label=(\alph*),leftmargin=*,itemsep=1pt]
\item \textbf{Prediction error}: The slow closure committed based on an incorrect model of the body or world. (Example: slipping on ice—the model predicted friction; the world lacked it.)
\item \textbf{Execution failure}: The fast closure lost its ability to close during execution. (Example: weak muscles, severed proprioception, sensory loss.)
\item \textbf{Unmodeled perturbation}: An external disturbance not represented in the slow closure's model disrupted execution. (Example: sudden gust of wind, unexpected obstacle.)
\end{enumerate}

In all cases, the outcome reveals that the slow closure did not intentionally commit to the harmful action. The action was not issued; the circuit failed.
\end{principle}

\begin{remark}[Deafferentation as circuit break]
A deafferented patient (complete loss of proprioceptive and large-fiber sensory input) cannot stand with eyes closed because the fast-closure proprioceptive return path is severed. The circuit cannot close. No bias signal from the slow closure can compensate, because the fast closure literally lacks the feedback to stabilize. The patient falls not because the slow closure committed to falling, but because the circuit topology is broken. Standing requires proprioceptive return; without it, the fast loop cannot function. The fall is a topological necessity, not a choice.
\end{remark}

\begin{remark}[F1 driver crash as prediction error]
If an elite F1 driver crashes into a barrier at a corner where they have safely navigated 100 times, the slow closure did not commit to crashing. The crash reveals either:
\begin{itemize}
\item Prediction error: model said the car would brake in time; the car did not (brake failure).
\item Execution failure: sensorimotor system lost control authority (steering failure, loss of grip).
\item Unmodeled perturbation: unexpected obstacle, mechanical failure, environmental change.
\end{itemize}

All reveal circuit failure or prediction error, not intentional commitment. The outcome (harm) cannot be desired; therefore, the circuit either predicted incorrectly or failed to execute.
\end{remark}

\subsection{Stability and Attractor Geometry of the Slow Closure}

\begin{definition}[Slow-closure attractor]
The slow closure, like the fast closure, operates near a limit cycle or fixed point in the space of internal representations. This attractor is defined by the integrated supraspinal dynamics: dopaminergic reward signals, prefrontal working memory, emotional state representations, and long-term goals.

The attractor can shift based on:
\begin{itemize}
\item Learning (synaptic weight changes from experience)
\item Motivation (dopaminergic state, limbic input, goal specification)
\item External constraints (instructions, task requirements)
\item Physiological state (fatigue, intoxication, metabolic demand)
\end{itemize}
\end{definition}

\begin{proposition}[Slow-closure commitment reflects attractor geometry]
The bias signals $b(t)$ issued by the slow closure form a trajectory on the attractor basin. Deliberation explores the basin (high variance, $\sigma > 0$). Commitment selects a point on the basin and executes it ($\sigma = 0$). The set of possible commitments is constrained by the attractor's geometry, not by arbitrary choice.
\end{proposition}

\subsection{Summary: Temporal Hierarchy as Substrate for Experience}

The temporal hierarchy of circuit closures—fast reflex, medium coordination, slow deliberation—provides the foundational topology for what is phenomenologically called ``conscious experience.''

Deliberative charge redistribution (slow closure) is the physical correlate of the subjective experience of planning, considering options, and deciding. Executive charge redistribution (fast closure) is the correlate of the subjective experience of acting and observing the consequences.

The hierarchy explains why:
\begin{itemize}
\item \textbf{Some actions are automatic}: The fast closure closes without slow-closure involvement. (Example: catching a falling object.)
\item \textbf{Some decisions feel difficult}: The slow closure explores a large basin with many competing attractors. (Example: choosing between careers.)
\item \textbf{Some actions feel natural}: The slow-closure commitment aligns with the current attractor. (Example: reaching for a familiar object.)
\item \textbf{Some actions feel forced}: The slow closure issues a commitment that conflicts with the emotional/motivational attractor. (Example: forcing yourself to study when tired.)
\end{itemize}

Pathological disruption at each level produces characteristic syndromes:
\begin{itemize}
\item \textbf{Fast-closure disruption} (spinal cord injury, peripheral neuropathy): Movement becomes impossible despite preserved slow-closure function. The patient ``wants to move'' but cannot.
\item \textbf{Slow-closure disruption} (prefrontal damage, severe depression): Movement remains mechanically possible, but goal-directed action is impaired. The patient is ``stuck in the moment'' without future planning.
\item \textbf{Coupling disruption} (cerebellar pathology, schizophrenia): Both closures function but lose phase coherence. Actions are executed but feel uncontrolled; deliberation loses grounding in executable actions.
\end{itemize}